\documentclass[11pt]{article}

\usepackage[margin=1in]{geometry}
\usepackage{amsmath,amssymb}
\usepackage{graphicx}
\usepackage{booktabs}
\usepackage{hyperref}
\usepackage{natbib}
\usepackage{enumitem}
\usepackage{xcolor}
\usepackage{caption}

\hypersetup{colorlinks=true, linkcolor=blue, citecolor=blue, urlcolor=blue}

\title{Testing and Correcting for Insufficient Follow-Up\\
in Covariate-Dependent Cure Models:\\
\large From Categorical Tests to a Heavy-Tailed Extrapolation, and Back}
\author{Anuj Kumar}
\date{August 2026}

\begin{document}
\maketitle

\begin{abstract}
Cure models are widely used to analyze time-to-event data in which a subpopulation
never experiences the event of interest: for example, cancer patients who are
effectively cured, or borrowers who will never default. A central practical
difficulty in fitting these models is \emph{insufficient follow-up}: if the observation
window is too short relative to the tail of the event-time distribution, the cured
fraction is systematically overestimated, because patients who would eventually
experience the event but have not yet done so are indistinguishable, within the
observed window, from those who are truly cured. Extreme value theory (EVT) offers
a principled way to relax the traditional ``sufficient follow-up'' assumption by
modeling the tail of the susceptible population's survival distribution directly.
This report documents an applied project built around four recent papers on this
problem, combining (i) a statistical test for whether follow-up is long enough,
(ii) an EVT-based correction to the cure-rate estimator when it is not, and
(iii) an investigation of how both tools behave as the number of covariates grows.
All methods are applied to the real \texttt{survival::colon} dataset (929
colon-cancer patients, two endpoints), rather than to synthetic data alone, and
several implementation bugs in both public reference packages and our own
from-scratch code were found and corrected along the way. The project was built as
a demonstration piece to accompany a PhD application to VU Amsterdam on exactly
this topic.
\end{abstract}

\tableofcontents

\section{Motivation and Problem Statement}

Standard survival analysis assumes that, given enough time, every subject in the
study population will eventually experience the event of interest. \emph{Cure
models} relax this by positing a latent split of the population into a
\emph{susceptible} fraction, who will eventually experience the event, and a
\emph{cured} (or immune) fraction, who never will. This structure is natural in
oncology (a meaningful fraction of treated cancer patients never relapse) and in
credit risk (a meaningful fraction of borrowers never default over the life of a
loan).

Estimating the cured fraction from data requires distinguishing, among subjects who
have not yet had the event by the end of the study, those who never will from those
who simply have not yet reached their (long) event time. This is only identifiable
if the study's follow-up is long enough that essentially all susceptible subjects
have already had the event by the end of observation: the \emph{sufficient
follow-up} (SFU) assumption. When follow-up is insufficient, naive nonparametric
estimators (e.g. Kaplan--Meier-based cure-rate estimates) are biased: they
systematically \emph{overestimate} the cured fraction, because long-surviving
susceptible subjects are counted as if cured.

The applied problem this project addresses has two parts:
\begin{enumerate}[nosep]
    \item \textbf{Detection.} Given a dataset, is follow-up actually long enough to trust
    a naive cure-rate estimate, overall and within covariate subgroups?
    \item \textbf{Correction.} If not, can the tail behavior of the event-time
    distribution be modeled well enough, using extreme value theory, to correct the
    naive estimate, and does this correction remain usable as more covariates
    are added?
\end{enumerate}

This project was built around four recent papers that address exactly these two
questions from different angles, and applies all of them, end to end, to a single
real dataset, rather than treating each paper as an isolated synthetic-data exercise.


\section{Literature Anchors}


Four papers anchor the project. All four were read in full and their key
derivations were worked through by hand before any code was written.

\begin{description}[leftmargin=1.2em, style=nextline]
    \item[Yuen \& Musta (2024/2026)] \citep{yuenmusta2024} An unconditional
    (marginal, no covariates) hypothesis test for sufficient follow-up, based on a
    smoothed Grenander-type estimator of the boundary of the event-time
    distribution's support. This is the baseline detection tool.

    \item[Yuen, Musta \& Van Keilegom (2025)] \citep{yuenmustavankeilegom2025}
    Extends the above test to categorical covariates. Proposes two complementary
    procedures: \emph{Method 1}, an intersection-union test (IUT) that only
    declares ``sufficient follow-up for all subgroups'' if every subgroup
    individually clears significance; and \emph{Method 2}, a bootstrap procedure
    that identifies which specific subgroup is most likely to have insufficient
    follow-up.

    \item[Xie, Escobar-Bach \& Van Keilegom (2023/2024)] \citep{xieebvk2023} A
    sufficient-follow-up test in the Gumbel maximum-domain-of-attraction case, using
    an Aitken-$\Delta^2$-style correction. Uses the opposite null-hypothesis
    convention to Yuen \& Musta ($H_0$: sufficient, rather than $H_0$: insufficient),
    which matters when combining results across papers.

    \item[Escobar-Bach \& Van Keilegom (2019/2023)] \citep{ebvk2023} The
    \emph{correction} paper, and the anchor for the extrapolation component of this
    project. In the Fr\'echet maximum-domain-of-attraction (heavy-tailed) case, it
    constructs a Beran-type conditional Kaplan--Meier estimator, then corrects its
    tail using a Hill-type extreme-value index estimator $\hat\gamma(x)$ and an
    extrapolated cure-probability estimator $\hat p(x)$, allowing the cure
    probability to be estimated \emph{beyond} the observed follow-up horizon.
\end{description}


\section{Data}


All methods are applied to \texttt{survival::colon}, a public dataset of 929
patients from a colon cancer clinical trial, with two competing endpoints:
\emph{recurrence} and \emph{death} (all-cause mortality). Covariates used across
the project include age, sex, treatment arm, number of positive lymph nodes
(\texttt{nodes}), tumor differentiation (\texttt{differ}), local extent of
invasion (\texttt{extent}), and indicators for obstruction, perforation, and
adherence to nearby organs.

Real data, rather than synthetic data alone, was used deliberately throughout: the
literature anchors mostly demonstrate their methods on synthetic examples or a
single toy application, and part of the point of this project is to see whether
the qualitative conclusions replicate, and where the methods break, on a genuine
dataset with realistic covariate structure, censoring, and sample sizes.

A key modeling choice is the cure/tail threshold $\tau$ (the point beyond which the
distribution's tail behavior is extrapolated). For recurrence, $\tau$ was set to
$\tau_G + 2$ years ($\approx 11.1$ years), justified by external clinical evidence:
a large pooled analysis (Pastorino et al., 2025, \textit{JAMA Oncology}, $n=35{,}213$
across 15 trials) shows colon-cancer recurrence risk falls below 0.5\% by six years
post-surgery. For death (all-cause mortality), no such disease-specific cure horizon
exists, so a more generous $\tau_G + 10$ years was used. This tau-sensitivity turned
out to matter substantively: an earlier, more naive placeholder tau (1.5$\times$ the
generic horizon) produced only a borderline recurrence result ($p \approx 0.06$–$0.08$),
while the clinically justified tau produced a clearly significant one ($p=0.023$),
illustrating that the choice of extrapolation horizon is not a nuisance parameter
but a substantive, defensible modeling decision.


\section{Phase 1: Baseline Replication of the Sufficient Follow-Up Tests}


\subsection{Method}

The unconditional and categorical sufficient-follow-up tests of
\citet{yuenmusta2024} and \citet{yuenmustavankeilegom2025} were applied to both
endpoints of the \texttt{colon} data, using the authors' own reference
implementation, the \texttt{cureSFUTest} R package. Two categorical splits were
tested: age (binned at the median) and sex.

\subsection{Implementation issues found and fixed}

Working with the real reference implementation surfaced concrete issues, all
resolved and documented rather than worked around silently:

\begin{itemize}[nosep]
    \item \texttt{sfu.cov.test()} silently requires a numeric covariate matrix and
    raises a misleading error message otherwise.
    \item The package's own bootstrap loop reads \texttt{.Random.seed} without
    checking it exists, which crashes on a fresh R session; fixed by explicitly
    calling \texttt{set.seed()} beforehand. This is a genuine bug in the published
    reference implementation, worth citing as such.
    \item \texttt{sfu.test()}'s \texttt{method} argument silently resolves to
    \texttt{"g"} (plain Grenander) rather than \texttt{"sg"} (smoothed Grenander)
    if left at its default, which produces a degenerate $p=0$ on this data;
    \texttt{method = "sg"} must be passed explicitly to obtain the
    smoothed-Grenander test used throughout the rest of this report. (By
    contrast, \texttt{sfu.cov.test()}, used for the categorical splits below,
    always routes internally through the smoothed-Grenander bandwidth regardless
    of arguments supplied, so it is not exposed to this particular default-argument
    trap.)
    \item The initial hand-written aggregation of \emph{Method 1} (the
    intersection-union test) both (a) used the wrong, unconditional bandwidth
    instead of the paper's undersmoothed conditional bandwidth, and (b) aggregated
    per-subgroup $p$-values with $\min(p_x)$ rather than the correct
    $\max(p_x)$ implied by the intersection-union logic ($H_0$: insufficient for
    \emph{some} $x$; $H_1$: sufficient for \emph{all} $x$). Both were fixed --
    the corrected implementation now reads Method 1's per-subgroup $p$-values
    directly off the same \texttt{sfu.cov.test()} call already used for Method 2,
    which internally always uses the correct undersmoothed bandwidth
    ($bw = \min(\tau n^{-7/30},\, \tau/2)$). Getting the aggregation rule backwards
    would have silently flipped the recurrence-by-sex conclusion, so this
    correction mattered.
\end{itemize}

\subsection{Results}

\begin{table}[h]
\centering
\caption{Unconditional sufficient-follow-up test results, by endpoint.}
\begin{tabular}{lcc}
\toprule
Endpoint & $p$-value & Conclusion \\
\midrule
Recurrence & 0.023 & Sufficient follow-up \\
Death (all-cause) & 0.66 & Insufficient follow-up \\
\bottomrule
\end{tabular}
\end{table}

\begin{table}[h]
\centering
\caption{Categorical (Method 1, corrected) sufficient-follow-up test, per-subgroup $p$-values.}
\begin{tabular}{lcc}
\toprule
Split & Subgroup $p$-values & Method 1 ($\max p$) conclusion \\
\midrule
Recurrence $\times$ sex & female 0.001, male 0.108 & Not sufficient for \emph{all} (male not significant) \\
Recurrence $\times$ age & 0.015, 0.043 & Sufficient for both age groups \\
Death $\times$ sex & 0.328, 0.122 & Insufficient \\
Death $\times$ age & 0.400, 0.138 & Insufficient \\
\bottomrule
\end{tabular}
\end{table}

The headline pattern is stable across every subgroup and survives the
bandwidth/aggregation bug fix unchanged: recurrence follow-up is adequate, while
death (all-cause mortality) follow-up is not. This makes sense clinically: colon-cancer
recurrence has a genuine, relatively short biological horizon, whereas all-cause
death has no such horizon within a finite trial follow-up window.

The recurrence-by-sex split is the most interesting finding of this phase: female
patients show clearly sufficient follow-up ($p=0.001$), while male
patients are borderline and not statistically distinguishable from insufficient
($p=0.108$). Age, by contrast, does not meaningfully split the recurrence
conclusion. This is exactly the kind of covariate-dependent effect that a purely
unconditional test would miss.

A secondary finding, useful methodologically: \citet{yuenmustavankeilegom2025}'s
own paper flags that Method 2's bootstrap-based worst-category selection can
disagree with a naive largest-$p$-value ranking (``selection instability''). This
was independently reproduced on the death-by-age and death-by-sex splits, across
two separate runs.


\section{Phase 2: Extrapolation Estimator and the Curse of Dimensionality}


\subsection{Method}

No existing R package implements the conditional, heavy-tailed correction
estimator of \citet{ebvk2023}, so it was implemented from equations (2.6)--(2.8) of
that paper. The base layer is a Beran (1981) kernel-weighted conditional
Kaplan--Meier estimator (via \texttt{condSURV::Beran}, Epanechnikov kernel,
bandwidth from \texttt{KernSmooth::dpik}). On top of this, a three-point ratio
estimator gives the local tail index $\hat\gamma(x)$, and a two-point estimator
gives the extrapolated cure probability $\hat p(x)$, using two geometrically-spaced
tuning constants $(y_1, y_2)$ selected over a grid via the paper's ``quadratic
errors criterion.'' This selection criterion was algebraically simplified to an
equivalent, cheaper form: choosing the grid point whose $\hat p$ is closest to the
mean of all grid $\hat p$ values, since
$\arg\min_k \sum_j (P_k - P_j)^2 \Leftrightarrow \arg\min_k |P_k - \text{mean}(P)|$.

\subsection{Validation and bugs found}

The implementation was validated against synthetic data with known ground truth
($\gamma(x)=(x{+}1)/2$, a logistic cure-probability function $p(x)$, Pareto-tailed
uncured survival times, and an artificially short follow-up window to force a
genuine insufficient-follow-up bias). The first run produced impossible values
($\hat p > 1$), traced to a missing implementation detail: the source paper
truncates $\hat\gamma(x)$ from below at 0.1 before it enters the $\hat p$ formula's
$y_1^{-1/\gamma}$ term. After adding this truncation, the estimator recovered
sensible values close to ground truth, and moved the naive Beran estimate in the
theoretically correct direction (upward, since naive Beran always underestimates
the true non-cure probability under insufficient follow-up).

Applying the validated estimator to the real \texttt{colon} data surfaced a second,
real-data-specific bug: on real data (a sparser effective local sample near the
tail than the synthetic validation setup), $\hat\gamma(x)$ can swing very large
(up to $\approx 40$ at age 75), which sends $\hat p(x)$ outside $[0,1]$ through the
denominator $y_1^{-1/\gamma}-1$ approaching zero. This was fixed by restricting the
tuning grid to candidates whose implied $\hat p$ falls in the theoretically required
range $[F_n(\tau_n), 1]$ before applying the closest-to-mean selection rule.

\subsection{Result: recurrence vs.\ death, and triangulation with Phase 1}

Applied to real \texttt{colon} data (age covariate, 9-point grid over ages 35--75),
the corrected estimator satisfies $\hat p(x) \geq$ naive Beran everywhere, as
theory requires, and the correction is markedly larger for death (mean gap 0.078)
than for recurrence (mean gap 0.038). This is a genuine triangulation: two
independent methods (the Phase 1 hypothesis test and the Phase 2 extrapolation
estimator) \emph{and} the source paper's own colon application all agree
that recurrence follow-up is close to adequate while death follow-up is not.

A sex-stratified check (splitting the same age-grid analysis by sex, to compare
against the Phase 1 recurrence-by-sex finding) did \emph{not} replicate cleanly:
the mean correction gap was larger for females (0.029) than males (0.021), the
opposite direction from the Phase 1 SFU test, and per-age gaps were noisy (e.g.\
one age point had no valid tuning pair at all). This is attributed to sample
thinning: halving the sample per stratum shrinks the local effective sample size
the Beran estimator sees near the tail, making the tail-index estimation
noticeably noisier. This non-replication is reported honestly rather than forced
into agreement with Phase 1; the pooled recurrence-vs-death result remains the
stronger, cleanly-triangulated finding.

\subsection{Curse-of-dimensionality demonstration}

The core methodological question the vacancy for this project is built around,
statistical learning for high-dimensional covariates with theoretical
guarantees, motivated a concrete, mechanism-level demonstration of \emph{why}
kernel-based conditional estimators like Beran's break down as covariates
accumulate, rather than a single abstract assertion that the curse of
dimensionality applies. The Beran estimator was extended to a multivariate
product-Epanechnikov kernel, and real \texttt{colon} recurrence covariates were
added one at a time (age, then $+$nodes, $+$differ, $+$extent), evaluated at the
componentwise-median ``typical patient,'' with a \emph{fixed} per-dimension
bandwidth (from \texttt{dpik}, or a Silverman-rule fallback for near-degenerate
ordinal covariates) reused unchanged as dimensions were added. Holding
bandwidths fixed as dimensions grow ensures that any degradation reflects
dimensionality itself, not bandwidth retuning.

\begin{table}[h]
\centering
\caption{Curse-of-dimensionality demonstration: effective sample size collapses as covariates accumulate.}
\begin{tabular}{lcccc}
\toprule
Covariates & Kish ESS & naive $p$ & corrected $\hat p$ & $\hat\gamma$ \\
\midrule
age & 134.8 & 0.508 & 0.546 & 0.936 \\
$+$ nodes & 36.6 & 0.359 & 0.395 & 1.214 \\
$+$ differ & 29.4 & 0.362 & 0.418 & 1.241 \\
$+$ extent & 18.4 & 0.495 & 0.496 & \textbf{0.100 (truncation floor)} \\
\bottomrule
\end{tabular}
\end{table}

The effective sample size (Kish's ESS) collapses monotonically from 135 to 18 as
four covariates are added with a fixed bandwidth. By four covariates, the
tail-index estimator has visibly broken down: $\hat\gamma$ is pinned exactly at the
0.1 truncation floor introduced during validation, a clean numerical signature of
outright estimation failure rather than a gradual, still-informative decline.
This is concrete, auditable evidence of the curse-of-dimensionality problem, not an
assertion.

A companion analysis multiplying dimensions differently (age, $+$nodes, $+$sex,
$+$differ, evaluated on the death endpoint with $n=888$) reinforced this at a
different evaluation point and additionally showed that the failure mode is not
uniform across covariate space. Repeating the correction-gap analysis across five
evaluation points (young/old age $\times$ low/high \texttt{nodes} profiles, plus
the median patient) revealed that the clean ``gap shrinks monotonically with
dimension'' story from a single median-point run does \emph{not} generalize: the
mean correction gap trends down through three added covariates but ticks back up
at the fourth, and $\hat\gamma$ floor-pinning is non-monotone across dimension.

The real driver turned out to be \emph{local covariate-value rarity}, not raw
dimension count: evaluation points with unusual \texttt{nodes} values (e.g.\
nodes$=8$, versus a typical range of 1--5) collapse to an effective sample size of
only 7--11 already at two covariates, and show $\hat\gamma$ bouncing erratically as
further dimensions are added. Central, typical points (the median patient,
low-nodes patients) stay comparatively stable (ESS 60--160) even at four
covariates.

A subsequent $B=150$ bootstrap check at three representative points (the median
patient, a young/high-nodes patient, and an old/high-nodes patient), repeated
across all four dimensions, confirmed this is a resampling-stable pattern rather
than a single-sample artifact: for the typical/median patient, $\hat\gamma$
floor-pinning rises from 5\% of bootstrap resamples at one covariate to
63--69\% at three to four covariates, a systematic, dimensionality-driven
collapse, while atypical high-nodes points are already unstable at a single
covariate and simply stay noisy rather than worsening further, a small-sample
fragility that dimensionality compounds rather than causes.

This yields a sharper conclusion than a single blanket curse-of-dimensionality
statement: there are \emph{two separable failure modes}: systematic, genuine
dimensionality collapse for typical patients, and pre-existing small-sample
fragility for atypical, rare covariate profiles that dimensionality merely
compounds.

\subsection{Comparison against a gradient-boosted alternative (gbex)}

To test whether a modern high-dimensional tail-regression tool resolves either
failure mode, \texttt{gbex} \citep{gbex2021} (a gradient-boosted generalized
Pareto regression method, co-authored by one of this project's target
supervisors) was installed from source (worked around a GitHub API rate limit by
downloading the tarball directly via \url{codeload.github.com} rather than
\texttt{remotes::in\-stall\_git\-hub}) along with its undeclared dependencies
(\texttt{POT}, \texttt{treeClust}).

Getting \texttt{gbex} to run correctly on this data required fixing several real
usage bugs: the function signature does not accept a \texttt{status}/
\texttt{threshold}/\texttt{maxit}/\texttt{stopcrit} interface (an initial
implementation attempt assumed it did); \texttt{predict()} does not support
\texttt{type="response"} and instead returns a data frame with columns \texttt{s}
(scale) and \texttt{g} (shape); covariates must be passed as a proper data
frame with real column names, not a bare matrix, or \texttt{predict.gbex} silently
mismatches columns internally; and, encountered later when computing variable
importance for the ten-covariate model below, \texttt{gbex::variable\_importance()}'s
default \texttt{type} argument (\texttt{type = c("relative","permutation")}) is
never resolved internally via \texttt{match.arg()}, so the default call errors
with ``the condition has length > 1'' -- \texttt{type} must be passed explicitly.

A more fundamental limitation was then found: \texttt{gbex} has no built-in notion
of censoring, so it can only be fit on \emph{uncensored} exceedances above a
threshold. Counting usable points across candidate thresholds showed this is
severely limiting on this dataset: for recurrence (whose follow-up is already
adequate per Phase 1, so there is little tail-exceedance signal left to model), the
usable count drops from 32 points at the median threshold to just 1 at the 75th/90th
percentile; for death, from 33 at the median threshold to 4 at the 90th percentile.
Given this, the comparison was deliberately scoped down to the death endpoint at a
low (50th-percentile) threshold, where 33 uncensored exceedances are available.

On this constrained but usable setup, \texttt{gbex} produced a tail-index estimate
$\hat\gamma \approx 0.005$ (near-exponential) that was essentially \emph{constant}
across every evaluation point and every covariate dimension tested. Rather than
accepting this as genuine stability, the flat result was investigated further. The
root cause: with only 33 observations, \texttt{gbex}'s internal unconditional
maximum-likelihood initial guess for $\gamma$ came back negative (nonsensical for a
generalized Pareto tail), silently triggering a fixed default initialization
($\gamma=0.01$); from that fixed starting point, the boosting shrinkage is tiny
($\lambda \approx 0.007$), so the reported estimate barely moves regardless of the
covariates supplied. A 30-repetition bootstrap confirmed the estimate is frozen
under resampling too (SD $=0.00012$), i.e.\ the estimator is not responding to the
data at all, not merely stable.

This is an important negative finding, arguably more important than a
clean positive result would have been: at this sample size, neither method
reliably recovers patient-level tail heterogeneity, but the two methods fail in
qualitatively different ways. The Beran-based estimator fails \emph{visibly} --
explicit truncation-floor pinning is an auditable, detectable failure signature.
\texttt{gbex} fails \emph{silently}: it produces what looks like a converged,
stable estimate, but the stability is an artifact of a fallback initialization
rather than genuine convergence. The silent failure mode is arguably the more
practically dangerous one, since it could mislead a practitioner who does not
inspect the initialization behavior.

Separately, \texttt{gbex} was also applied to the recurrence endpoint using all ten
available real covariates simultaneously (age, sex, obstruction, perforation,
adherence, differentiation, extent, surgery type, nodes, and a nodes$\geq$4
indicator) on 223 uncensored exceedances, explicitly as a demonstration of
covariate-scaling for the \emph{scale} parameter alone, not a claim of solving the
censoring/cure problem directly (this remains the open methodological gap the
target PhD project is meant to address). Here \texttt{gbex} performed
well: the estimated scale parameter varied substantially and meaningfully across
patients (0.31--3.13), driven mainly by age, number of positive nodes,
differentiation, and obstruction, and the full ten-covariate model outperformed an
age-only model on deviance (1.221 vs.\ 1.280). The estimated shape (tail index),
however, stayed pinned near zero for every patient regardless of covariates --
consistent with the unconditional estimate, which was actually negative before
\texttt{gbex}'s positivity constraint clamped it. This suggests that, among
patients who do recur, the raw recurrence-time distribution shows little evidence
of genuine heavy-tailedness; the insufficient-follow-up problem in this data more
plausibly comes from the cure/mixture structure itself (a real never-recur
subpopulation) than from a heavy-tailed susceptible population. This is a useful
methodological point in its own right: it argues for the mixture-cure framing used
throughout this project over a naive heavy-tail-only model, while also
demonstrating \texttt{gbex}'s genuine value on the scale parameter specifically.


\section{Phase 5: Simulation Study}

Phase 5 was prioritized over the originally planned Phase 3 (a continuous-covariate
sufficient-follow-up test) and the full Phase 4 (a comprehensive per-estimator
uncertainty-quantification framework) because it reused already-validated code at
low marginal cost and produced results with direct interpretive value. Phase 3 was
explicitly scoped as a stretch or sketch item from the outset, and the full Phase 4
was judged to be substantial, open-ended work relative to the remaining time before
the application deadline.

\subsection{Part A: Bias--variance tradeoff as follow-up length varies}

Using the validated Phase 2 estimator on $N=2000$ synthetic patients (true
$p(x_0=0.5) = 0.5987$), the censoring horizon $\tau_c$ was swept over
$\{4, 6, 8, 12, 20, 40\}$, 30 replications per value.

\begin{table}[h]
\centering
\caption{Bias and MSE of naive Beran vs.\ EBVK-corrected estimator as follow-up lengthens ($\tau_c$).}
\begin{tabular}{lcccc}
\toprule
$\tau_c$ & Bias (naive) & Bias (corrected) & MSE (naive) & MSE (corrected) \\
\midrule
4 (severe) & $-0.095$ & $-0.038$ & 0.0119 & \textbf{0.0038} \\
6 & $-0.038$ & $+0.018$ & 0.0044 & 0.0039 \\
8 & $-0.031$ & $+0.012$ & 0.0042 & 0.0053 \\
12 & $-0.024$ & $-0.003$ & 0.0027 & 0.0027 \\
20 & $-0.002$ & $+0.009$ & 0.0015 & 0.0018 \\
40 (near-sufficient) & $+0.003$ & $+0.011$ & 0.0022 & 0.0026 \\
\bottomrule
\end{tabular}
\end{table}

This demonstrates the bias-variance tradeoff the correction is theoretically
expected to exhibit, though the pattern is not perfectly monotonic. At severe
insufficiency ($\tau_c=4$), the corrected estimator more than halves the bias and
wins decisively on MSE (0.0038 vs.\ 0.0119), exactly where correcting matters
most, and it still wins narrowly at $\tau_c=6$ (0.0039 vs.\ 0.0044). By $\tau_c=8$
the correction's MSE advantage has already flipped (0.0053 vs.\ 0.0042), before
the two estimators tie again at $\tau_c=12$ and naive Beran pulls ahead from
$\tau_c=20$ onward. With only 30 replications per $\tau_c$, the dip at $\tau_c=8$
is plausibly Monte Carlo noise rather than a genuine non-monotonicity, and it
does not change the overall picture: as follow-up approaches sufficiency, naive
Beran's own bias shrinks toward zero on its own, so the correction's advantage
fades and eventually reverses, since the correction adds variance without a
meaningful bias left to remove. This confirms, on a controlled synthetic sweep
rather than a single point estimate, that the correction should be applied when
follow-up is insufficient and can be safely skipped once it is close to
adequate.

\subsection{Part B: Method 1 vs.\ Method 2 as the number of covariate categories grows}

Using the real \texttt{cureSFUTest} package, the number of covariate categories
$K$ was swept over $2$--$5$, with a fixed total sample size $n=400$ split evenly
across categories (so per-category $n$ shrinks with $K$: 200, 133, 100, 80).

An initial pass mislabeled which simulated scenario corresponds to which error
type; this was caught and corrected before reporting results below. In the
``all categories sufficient'' scenario the ground truth is $H_1$, so \emph{failing}
to detect it is a Type~II error, and the complement is the test's power to confirm
true full sufficiency. In the ``one category insufficient'' scenario the ground
truth is $H_0$, so wrongly declaring ``all sufficient'' is the actual Type~I error.

\begin{table}[h]
\centering
\caption{Method 1 (intersection-union test) power and Type~I error, and Method 2 selection accuracy, as $K$ grows.}
\resizebox{\textwidth}{!}{%
\begin{tabular}{lcccc}
\toprule
$K$ & $n$/category & Power (confirm true sufficiency) & Type I error (false ``all clear'') & Method 2 selection accuracy \\
\midrule
2 & 200 & 0.60 & 0.067 & 1.00 \\
3 & 133 & 0.43 & 0.00 & 0.967 \\
4 & 100 & 0.27 & 0.00 & 0.967 \\
5 & 80 & 0.30 & 0.00 & 1.00 \\
\bottomrule
\end{tabular}%
}
\end{table}

This produces a complementary-strengths finding. Method 1's Type~I error
(falsely declaring ``sufficient follow-up for all subgroups'' when at least one
subgroup truly is not) stays at or near zero across every $K$ tested, and if
anything the test becomes \emph{more} conservative as $K$ grows: the
intersection-union test essentially never gives false reassurance. This safety has
a real cost, however: power to positively \emph{confirm} genuine full sufficiency
collapses from 60\% at $K=2$ to roughly 27--30\% at $K \geq 4$, because every one of
$K$ increasingly noisy, shrinking-sample-size, per-category tests must
simultaneously clear the 5\% significance bar. Method 2's bootstrap-based
diagnostic selection of \emph{which} category is problematic, by contrast, stays
reliable (97--100\% accurate) throughout, regardless of $K$. The practical
implication: Method 1 should be trusted for its conservatism (it will not falsely
clear a real problem) but not relied upon to certify ``everything is fine'' once
several covariate categories are involved; Method 2 remains the right tool for
localizing where a problem is once one is suspected.


\section{Discussion: What This Project Demonstrates}
Across all five completed phases, three findings stand out as the strongest
material for further development:

\begin{enumerate}[nosep]
    \item \textbf{Triangulated main finding (Phases 1 and 2).} Two independent
    methods, a hypothesis test and an extrapolation estimator, and the original
    source paper's own real-data application all agree: colon-cancer recurrence
    follow-up in this dataset is close to sufficient, while all-cause death
    follow-up is not. This agreement across independently implemented methods is
    strong evidence that the qualitative conclusion is real, not an artifact of one
    estimator.

    \item \textbf{Mechanism-level curse-of-dimensionality evidence (Phase 2).} Rather
    than asserting that kernel-based conditional cure estimators degrade with more
    covariates, this project traces the exact numerical mechanism: effective sample
    size collapse, and a clean truncation-floor signature marking the point of
    estimator breakdown. It further separates this into two distinct failure
    modes, systematic dimensionality collapse for typical patients, and
    pre-existing small-sample fragility for atypical ones, and shows that a
    modern gradient-boosted alternative (\texttt{gbex}) does not straightforwardly
    solve the problem: at low sample sizes it fails \emph{silently}, via a fallback
    initialization, which is arguably more dangerous in practice than the
    Beran estimator's visible truncation-floor failure.

    \item \textbf{Bias--variance tradeoff (Phase 5, Part A).} A controlled
    synthetic sweep over follow-up length reproduces, with a full multi-point
    curve rather than a single spot check, the theoretical bias--variance
    tradeoff the correction method is designed to exhibit: large, decisive gains
    when follow-up is severely insufficient, and a cost in variance once
    follow-up is already close to adequate, with one likely-noise dip at
    $\tau_c=8$ noted rather than smoothed over.
\end{enumerate}

An honest secondary finding worth naming explicitly is the sex-stratified
non-replication in Phase 2 (Section 4.3): the Phase 1 hypothesis test and the
Phase 2 extrapolation estimator disagree on the direction of the recurrence-by-sex
effect once the sample is split by sex, most plausibly because of sample-size
thinning rather than a genuine methodological contradiction. This is reported as
an honest limitation rather than adjusted to agree with the headline result.

Three concrete implementation bugs were found and corrected in the public
\texttt{cureSFUTest} reference package (Section 4.2): the numeric-covariate-matrix
requirement of \texttt{sfu.cov.test()}, the \texttt{.Random.seed} bootstrap crash,
and \texttt{sfu.test()}'s silent default-argument fallback to the unsmoothed
Grenander test. A further four usage bugs were found while fitting and evaluating
\texttt{gbex} (Section 4.4): the assumed \texttt{status}/\texttt{threshold}/
\texttt{maxit}/\texttt{stopcrit} interface, \texttt{predict()}'s unsupported
\texttt{type="response"} argument, the data-frame-vs-matrix column-mangling issue,
and \texttt{variable\_importance()}'s unresolved default \texttt{type} argument --
alongside the separate, more consequential silent-initialization failure mode
also documented in that section. Four more bugs were found and fixed in this
project's own from-scratch code: a missing tail-index truncation, an
out-of-range tuning-grid selection, an intersection-union aggregation-rule error,
and an incorrect Type~I/II error labeling in simulation code. All are documented
above rather than silently fixed, since correctly diagnosing and fixing
implementation bugs in both published reference packages and one's own code is
itself part of the applied statistical skill this project is meant to demonstrate.


\section{Future Work}

Two components of the original six-phase plan were deliberately deferred, and are
named here as concrete open problems rather than treated as gaps in the current
work.

\paragraph{A continuous-covariate sufficient-follow-up test.} The categorical test
of \citet{yuenmustavankeilegom2025} requires discretizing continuous covariates
(as was done here for age) into bins, which discards information and makes the
Phase 1 power problem (Section 6, Part B) worse as more categories are used. A
natural extension combines Yuen \& Musta's smoothed-Grenander boundary estimator
with the Beran-type conditional weighting used in Phase 2, replacing discrete
subsampling by continuous covariate weighting. This does not require a full
asymptotic theory to be useful as a first step: a working numerical version, paired
with an honest statement of exactly which asymptotic argument (e.g.\ the limiting
distribution of the smoothed-Grenander statistic under continuous local weighting)
remains open, is itself a meaningful contribution and a natural bridge into
doctoral-level theoretical work on this problem.

\paragraph{Per-estimator uncertainty quantification, done correctly.} The tests
used in Phases 1 and 5(B) come with their own bootstrap-based $p$-values, but the
Phase 2 extrapolation estimator ($\hat p(x)$, $\hat\gamma(x)$) currently has no
confidence intervals at all. A methodologically correct treatment requires
different resampling schemes for different estimator types: a naive bootstrap is
appropriate for smooth functionals (such as the Xie et al.\ and Escobar-Bach \&
Van Keilegom estimators), while a smoothed bootstrap is required for the
shape-constrained, Grenander-based tests, whose limiting distributions are
non-standard. Building this properly, per estimator type, is substantial
methodological work in its own right and is left as future work rather than
attempted with a one-size-fits-all interval that would likely be invalid for at
least one of the estimator types involved.

\paragraph{Resolving the silent gbex failure mode.} The Phase 2 comparison
surfaced a specific, fixable issue: \texttt{gbex}'s reliance on an unconditional
MLE initial guess for the tail index becomes unreliable at small uncensored sample
sizes, silently falling back to a fixed default. A natural methodological
extension, and a direct answer to the vacancy's stated interest in
high-dimensional covariates with theoretical guarantees, is to extend
\texttt{gbex} or a similar gradient-boosted tail-regression framework to handle
censoring natively (e.g.\ via a censored generalized Pareto likelihood), rather
than restricting it to uncensored exceedances as done here. This would substantially
increase the usable sample size for the recurrence endpoint in particular, where
follow-up is already adequate and censored near-threshold observations currently
carry no weight in the tail-regression step.


\section*{Acknowledgments}

This project was developed as a demonstration piece to accompany a PhD application
to the Econometrics and Data Science Department, School of Business and Economics,
VU Amsterdam, for the NWO-funded project \emph{``Seeing beyond the study duration:
extreme value theory for long-term survival and cure chances,''} supervised by
Dr.\ Juan Cai (VU Amsterdam) and Dr.\ Eni Musta (University of Amsterdam). All four
anchor papers were selected specifically because they originate from the target
supervisors' own research groups.


\begin{thebibliography}{9}

\bibitem[Yuen and Musta(2024)]{yuenmusta2024}
Yuen, T.\ and Musta, E.\ (2024/2026).
\textit{Testing for sufficient follow-up in survival data with a cure fraction.}
Biometrical Journal. arXiv:2403.16832.

\bibitem[Yuen et~al.(2025)]{yuenmustavankeilegom2025}
Yuen, T., Musta, E.\ and Van Keilegom, I.\ (2025).
\textit{Testing for sufficient follow-up in cure models with categorical covariates.}
arXiv:2505.13128.

\bibitem[Xie et~al.(2023)]{xieebvk2023}
Xie, R., Escobar-Bach, M.\ and Van Keilegom, I.\ (2023/2024).
\textit{Testing for sufficient follow-up in censored survival data by using extremes.}
Biometrical Journal, 66(4):e202400033. arXiv:2309.00868.

\bibitem[Escobar-Bach and Van Keilegom(2019)]{ebvk2023}
Escobar-Bach, M.\ and Van Keilegom, I.\ (2019/2023).
\textit{Nonparametric estimation of conditional cure models for heavy-tailed distributions and under insufficient follow-up.}
Computational Statistics \& Data Analysis, 183:107728. arXiv:1909.08436.

\bibitem[Velthoen et~al.(2021)]{gbex2021}
Velthoen, J., Dombry, C., Cai, J.-J.\ and Engelke, S.\ (2021/2023).
\textit{Gradient boosting for extreme quantile regression.}
Extremes, 26(4):639--667. arXiv:2103.00808.

\bibitem[Pastorino et~al.(2025)]{pastorino2025}
Pastorino, A.\ et al.\ (2025).
\textit{Timing and pattern of colon cancer recurrence: a pooled analysis of 15 adjuvant trials.}
JAMA Oncology.

\bibitem[Beran(1981)]{beran1981}
Beran, R.\ (1981).
\textit{Nonparametric regression with randomly censored survival data.}
Technical Report, University of California, Berkeley.

\end{thebibliography}

\end{document}
